\section{Square-zero deformation}
\label{sec:square_zero_deformation}

Let $A$ be a commutative ring, let $I\subseteq A$ be an ideal with
$I^2=0$, and put $B=A/I$.  In this section we prove that a real Novikov
descent datum over $A$ is constant whenever its reduction to $B$ is constant.
The argument has two inputs: deformation theory for finite projective modules
and an additive \v{C}ech calculation for Novikov series.

\subsection{Square-zero algebra for finite projective modules}

\begin{lemma}[Maps vanishing after square-zero base change]
  \label{lem:square_zero_base_change}
  \lean{Novikov.Miscellany.reductionMap,
    Novikov.Miscellany.reductionMap_surjective,
    Novikov.Miscellany.reductionMap_ker,
    Novikov.Miscellany.comp_eq_zero_of_baseChangeMap_eq_zero}
  \leanok
  Let $q\colon A\twoheadrightarrow B$ be a surjective homomorphism with
  $K=\ker(q)$ and $K^2=0$.  For every $A$-module $M$, the reduction map
  \[
    M\longrightarrow B\otimes_A M,\qquad m\longmapsto 1\otimes m,
  \]
  is surjective and has kernel $KM$.  Consequently, if composable
  $A$-linear maps $v\colon L\to M$ and $u\colon M\to N$ both become zero
  after base change to $B$, then $u\circ v=0$.
\end{lemma}

\begin{proof}
  \leanok
  Surjectivity follows by writing a pure tensor $b\otimes m$ as
  $1\otimes am$ after choosing $a\in A$ with $q(a)=b$.  The kernel statement
  is obtained by factoring $q$ through $A/K$ and using the canonical
  identification
  \[
    (A/K)\otimes_A M\cong M/KM.
  \]
  If $u$ becomes zero, its image is contained in $KN$; similarly, the image
  of $v$ is contained in $KM$.  Hence the image of $u\circ v$ is contained in
  $K^2N=0$.
\end{proof}

\begin{lemma}[Lemma~4.6: finite projectives over a square-zero thickening]
  \label{lem:vect_square_zero_deform}
  \lean{Novikov.Miscellany.exists_linearEquiv_lift_of_surjective_of_ker_sq,
    Novikov.Miscellany.FiniteProjectiveModule.exists_lift_of_surjective_of_ker_sq,
    Novikov.Miscellany.exists_linearEquiv_lift_of_quotient_sq,
    Novikov.Miscellany.FiniteProjectiveModule.exists_lift_of_quotient_sq}
  \leanok
  \uses{lem:square_zero_base_change}
  Let $I\subseteq A$ satisfy $I^2=0$.
  \begin{enumerate}
    \item If $M$ and $N$ are finite projective $A$-modules, every
      $A/I$-linear isomorphism
      \[
        (A/I)\otimes_A M\xrightarrow{\sim}(A/I)\otimes_A N
      \]
      lifts to an $A$-linear isomorphism $M\xrightarrow{\sim}N$ whose base
      change is the given map.
    \item Every finite projective $A/I$-module $\overline M$ is isomorphic to
      $(A/I)\otimes_A M$ for some finite projective $A$-module $M$.
  \end{enumerate}
  More generally, both assertions hold for any surjection $A\twoheadrightarrow
  B$ with square-zero kernel.
\end{lemma}

\begin{proof}
  \leanok
  For (1), projectivity lifts the given map and its inverse to maps
  $f\colon M\to N$ and $g\colon N\to M$.  Set
  \[
    a=gf-1_M,\qquad b=fg-1_N.
  \]
  Both $a$ and $b$ vanish after base change.  By
  Lemma~\ref{lem:square_zero_base_change}, $a^2=b^2=0$.  Moreover
  $fa=bf$.  Replacing $g$ by
  \[
    g'=(1_M-a)g
  \]
  gives $g'f=1_M$ and $fg'=1_N$, so $f$ is the required lift.

  For (2), present $\overline M$ as the image of an idempotent matrix over
  $A/I$.  The entrywise map on matrix rings is surjective and its kernel is
  nilpotent, so the idempotent lifts to an idempotent matrix over $A$.  The
  image of the corresponding endomorphism of a finite free $A$-module is
  finite projective, and its reduction is $\overline M$.
\end{proof}

\subsection{Additive Novikov \v{C}ech cocycles}

\begin{lemma}[Square-zero coefficient maps for Novikov rings]
  \label{lem:novikov_square_zero_coefficients}
  \lean{Novikov.mapRingHom_surjective, Novikov.mapRingHom_ker_sq}
  \leanok
  \uses{def:novikov_series, thm:novikov_ring}
  Let $q\colon A\twoheadrightarrow B$ be surjective with square-zero kernel.
  For every finite set of variables, the coefficientwise homomorphism
  \[
    A((t_1^{\mathbb R},\ldots,t_r^{\mathbb R}))
      \longrightarrow
    B((t_1^{\mathbb R},\ldots,t_r^{\mathbb R}))
  \]
  is surjective and has square-zero kernel.
\end{lemma}

\begin{proof}
  \leanok
  Lift each nonzero coefficient of a series over $B$ to $A$; using zero away
  from the original support preserves the Novikov finiteness condition.  Thus
  the coefficientwise map is surjective.  A series in its kernel has every
  coefficient in $\ker(q)$.  Every coefficient of a product of two such
  series is a finite sum of products of elements of $\ker(q)$, and therefore
  vanishes.
\end{proof}

\begin{lemma}[Additive Novikov \v{C}ech cocycles are coboundaries]
  \label{lem:additive_novikov_cech_coboundary}
  \lean{Novikov.Descent.firstAxis,
    Novikov.Descent.firstAxis_apply,
    Novikov.Descent.additive_cocycle_eq_coboundary}
  \leanok
  \uses{def:novikov_substitute, thm:novikov_substitute_functorial}
  Let $H$ be an abelian group and let
  $x\in H((t^{\mathbb R},u^{\mathbb R}))$ satisfy
  \[
    \pi_{13}^*x=\pi_{23}^*x+\pi_{12}^*x.
  \]
  Define the first-axis restriction $y\in H((t^{\mathbb R}))$ by
  $y_d=x_{(d,0)}$.  Then
  \[
    x=\pi_1^*y-\pi_2^*y.
  \]
\end{lemma}

\begin{proof}
  \leanok
  Evaluate the cocycle equation coefficientwise.  At an exponent $(a,0,b)$
  with $a,b\neq0$, it shows that the coefficient $x_{(a,b)}$ vanishes.  At
  the origin it gives $x_{(0,0)}=x_{(0,0)}+x_{(0,0)}$, hence
  $x_{(0,0)}=0$ without any assumption on the characteristic.  Finally,
  evaluating at $(0,d,0)$ shows
  $x_{(0,d)}=-x_{(d,0)}$ for $d\neq0$.  These three calculations are exactly
  the displayed coboundary formula.
\end{proof}

\begin{lemma}[Endomorphism-valued additive cocycles]
  \label{lem:endomorphism_novikov_cech_coboundary}
  \lean{Novikov.Descent.exists_end_eq_sub_of_is_add_cocycle,
    Novikov.Descent.exists_end_eq_sub_of_is_add_cocycle_of_map_eq_zero}
  \leanok
  \uses{lem:additive_novikov_cech_coboundary, thm:novikov_base_change}
  Let $P$ be a finite projective $A$-module and let $C_r(P)$ denote its scalar
  extension to the $r$-variable real Novikov ring.  If
  \[
    \delta\in\operatorname{End}(C_2(P)),\qquad
    \pi_{13}^*\delta=\pi_{23}^*\delta+\pi_{12}^*\delta,
  \]
  then there exists $\eta\in\operatorname{End}(C_1(P))$ such that
  \[
    \delta=\pi_1^*\eta-\pi_2^*\eta.
  \]
  If coefficient extension along $q\colon A\to B$ sends $\delta$ to zero,
  then $\eta$ may be chosen to become zero after the same coefficient
  extension.
\end{lemma}

\begin{proof}
  \leanok
  Finite projectivity gives compatible identifications
  \[
    \operatorname{End}_{R_r}(R_r\otimes_A P)
      \cong R_r\otimes_A\operatorname{End}_A(P)
      \cong \operatorname{End}_A(P)((t_1^{\mathbb R},\ldots,t_r^{\mathbb R})).
  \]
  Under these identifications, normalized pullback along a face is the
  corresponding substitution.  Apply
  Lemma~\ref{lem:additive_novikov_cech_coboundary} with
  $H=\operatorname{End}_A(P)$ to obtain $\eta$.  The construction
  is first-axis restriction, which commutes with coefficient maps; hence if
  $\delta$ becomes zero, the chosen $\eta$ does as well.
\end{proof}

\subsection{Lifting a constant descent datum}

\begin{lemma}[Normalization of a constant underlying module]
  \label{lem:normalized_constant_descent}
  \lean{Novikov.Descent.normalizedPhi,
    Novikov.Descent.normalizedPhi_constant,
    Novikov.Descent.normalizedPhi_cocycle,
    Novikov.Descent.vectToNovikovDescent_baseChangeIso}
  \leanok
  \uses{def:novikov_descent_datum, def:constant_descent_functor,
    thm:novikov_base_change}
  Let $P$ be a finite projective $A$-module.  The two pullbacks of the
  underlying module of $\operatorname{Const}(P)$ admit canonical
  identifications with
  \[
    C_2(P)=A((t^{\mathbb R},u^{\mathbb R}))\otimes_A P.
  \]
  Under these identifications, a descent isomorphism becomes an endomorphism
  $\psi$ of $C_2(P)$; the constant descent isomorphism becomes the identity,
  and the multiplicative cocycle condition becomes
  \[
    \pi_{23}^*\psi\circ\pi_{12}^*\psi=\pi_{13}^*\psi.
  \]
  These identifications are compatible with extension of the coefficient ring.
\end{lemma}

\begin{proof}
  \leanok
  The identifications are the associativity and cancellation isomorphisms for
  the iterated tensor products defining constant descent.  Conjugating by
  them sends the constant descent map to the identity.  Their level-three
  associativity tetrahedra identify all three pullbacks coherently, so
  conjugation transports the original cocycle equation to the displayed
  equation.  The same tensor associativity maps commute with coefficient
  extension, which gives the final assertion.
\end{proof}

\begin{lemma}[Lemma~4.7: square-zero deformation of Novikov descent]
  \label{lem:novikov_descent_square_zero}
  \lean{Novikov.Descent.novikovDescent_squareZero}
  \leanok
  \uses{lem:vect_square_zero_deform,
    lem:square_zero_base_change,
    lem:novikov_square_zero_coefficients,
    lem:endomorphism_novikov_cech_coboundary,
    lem:normalized_constant_descent}
  Let $A$ be a commutative ring and let $I\subseteq A$ satisfy $I^2=0$.
  Let $M$ be a real Novikov descent datum over $A$.  If its base change to
  $A/I$ is isomorphic to a constant descent datum, then $M$ is itself
  isomorphic to a constant descent datum over $A$.
\end{lemma}

\begin{proof}
  \leanok
  Put $B=A/I$.  By hypothesis there are a finite projective $B$-module
  $\overline P$ and an isomorphism
  \[
    \operatorname{Const}(\overline P)\xrightarrow{\sim}M_B.
  \]
  Lemma~\ref{lem:vect_square_zero_deform} lifts $\overline P$ to a finite
  projective $A$-module $P$.  Compatibility of constant descent with
  coefficient extension then gives an isomorphism
  \[
    M_B\xrightarrow{\sim}\operatorname{Const}(P)_B.
  \]

  At level one, Lemma~\ref{lem:novikov_square_zero_coefficients} says that the
  map of Novikov rings is again a surjective square-zero extension.  Applying
  Lemma~\ref{lem:vect_square_zero_deform} lifts the underlying linear
  isomorphism to
  \[
    \theta\colon M\xrightarrow{\sim}\operatorname{Const}(P)
  \]
  on the underlying modules.  Transport the descent datum along $\theta$ and
  normalize its descent isomorphism as in
  Lemma~\ref{lem:normalized_constant_descent}; call the resulting
  endomorphism $\psi$.  Its reduction is the identity.  Set
  $\delta=\psi-1$.

  The normalized multiplicative cocycle reads
  \[
    (1+\pi_{23}^*\delta)(1+\pi_{12}^*\delta)
      =1+\pi_{13}^*\delta.
  \]
  Both face pullbacks of $\delta$ become zero over $B$.  Their composite is
  therefore zero by Lemma~\ref{lem:square_zero_base_change}, now at level
  three.  Expanding the preceding equation gives the additive cocycle
  \[
    \pi_{13}^*\delta
      =\pi_{23}^*\delta+\pi_{12}^*\delta.
  \]
  Lemma~\ref{lem:endomorphism_novikov_cech_coboundary} supplies
  $\eta\in\operatorname{End}(C_1(P))$, still zero after reduction, with
  \[
    \delta=\pi_1^*\eta-\pi_2^*\eta.
  \]

  Again by Lemma~\ref{lem:square_zero_base_change}, $\eta^2=0$ and
  $\pi_2^*\eta\circ\pi_1^*\eta=0$.  Thus
  \[
    \xi=1+\eta,
    \qquad
    \xi^{-1}=1-\eta,
  \]
  and
  \[
    \psi
      =1+\pi_1^*\eta-\pi_2^*\eta
      =(1-\pi_2^*\eta)(1+\pi_1^*\eta)
      =\pi_2^*(\xi^{-1})\circ\pi_1^*(\xi).
  \]
  Undoing the normalized tensor identifications says precisely that $\xi$ is
  an isomorphism from the transported datum to $\operatorname{Const}(P)$.
  Composing it with the transport isomorphism and taking the inverse yields
  $\operatorname{Const}(P)\xrightarrow{\sim}M$.
\end{proof}
